% ATTEMPT_STATUS: unresolved
% ATTEMPT_ORIGIN: new research, 2026-10-01
% TOP_PROBLEM: {"id": 1372, "title": "Tammes Problem", "category_id": 6, "category": {"id": 6, "name": "geometry", "display_name": "Geometry", "description": "Euclidean and non-Euclidean geometry, geometric structures.", "slug": "geometry", "order_index": 6, "created_at": "2026-07-31T15:26:25.671Z"}, "status": "open"}
% FIRST_POSTED: 2026-10-01
% Imported from: unsolvedmath_additional_500.tex; UnsolvedMath ID 1372
% !TeX program = lualatex
% Compile twice: lualatex 1372.tex
% Self-contained: no external bibliography, images, or \input files.
% Numeric section labels and visible numbers are original UnsolvedMath IDs.
\documentclass[11pt,a4paper]{article}
\usepackage[margin=24mm,headheight=14pt]{geometry}
\usepackage{fontspec}
\setmainfont{Latin Modern Roman}
\setsansfont{Latin Modern Sans}
\setmonofont{Latin Modern Mono}
\usepackage{amsmath,amssymb,mathtools}
\usepackage{unicode-math}
\setmathfont{Latin Modern Math}
\usepackage{microtype}
\usepackage{enumitem,xurl,needspace}
\usepackage[dvipsnames]{xcolor}
\usepackage{hyperref}
\hypersetup{unicode=true,colorlinks=true,linkcolor=MidnightBlue,urlcolor=MidnightBlue,
 pdftitle={UnsolvedMath Problem 1372},pdfauthor={Source-annotated compilation},bookmarksnumbered=true}
\usepackage{bookmark,tocloft,titlesec,fancyhdr}
\setlength{\cftsecnumwidth}{4.2em}
\cftsetpnumwidth{2.5em}
\cftsetrmarg{4em}
\titleformat{\section}{\Large\bfseries\raggedright}{\thesection}{0.7em}{}
\pagestyle{fancy}\fancyhf{}
\fancyhead[L]{\small\sffamily UnsolvedMath research attempt}
\fancyhead[R]{\small Original catalogue IDs}
\fancyfoot[C]{\thepage}
\setlength{\parindent}{0pt}
\setlength{\parskip}{5pt plus 1pt}
\setlength{\emergencystretch}{3em}
\setcounter{tocdepth}{1}\setcounter{secnumdepth}{1}
\setlist[itemize]{leftmargin=3em,itemsep=4pt,topsep=4pt,parsep=0pt}
\allowdisplaybreaks
\newcommand{\N}{\mathbb{N}}
\newcommand{\Z}{\mathbb{Z}}
\newcommand{\Q}{\mathbb{Q}}
\newcommand{\R}{\mathbb{R}}
\newcommand{\C}{\mathbb{C}}
\newcommand{\F}{\mathbb{F}}
\newcommand{\A}{\mathbb{A}}
\newcommand{\PP}{\mathbb{P}}
\newcommand{\E}{\mathbb{E}}
\newcommand{\eps}{\varepsilon}
\newcommand{\dd}{\,\mathrm d}
\newcommand{\sref}[2]{\hyperlink{source:#1:#2}{[S#2]}}
\newif\ifshowcatalogue
\showcataloguetrue
\newcommand{\cataloguescope}[1]{\ifshowcatalogue
 \paragraph{Statement recorded in the supplied source.}
 #1\fi}
\begin{document}
% UnsolvedMath ID 1372; source catalogue code GEOM-013

\setcounter{section}{1371}

\section{The Tammes Spherical Separation Problem}\label{1372}

{\small\sffamily UnsolvedMath ID 1372 · Geometry}

\subsection{Definitions and mathematical statement}

\paragraph{Shared notation.}Unless a section says otherwise, $\N=\{1,2,\ldots\}$,
$\Z$ is the ring of integers, $\Q,\R,\C$ are the rational, real, and complex
fields, $[N]=\{1,\ldots,\lfloor N\rfloor\}$, and $\log$ is the natural
logarithm. For real functions of a parameter tending to infinity,
$f\ll g$ and $f=O(g)$ mean $|f|\leq Cg$ eventually for some fixed $C$;
$f\gg g$ reverses this comparison; $f=o(g)$ means $f/g\to0$;
$f\sim g$ means $f/g\to1$; and $f\asymp g$ means both order bounds.
A subscript on $O$ or $\ll$ specifies allowed constant dependence.
The expression $n^{o(1)}$ in an upper bound means growth smaller than
$n^\epsilon$ for each fixed $\epsilon>0$. Local definitions govern any
exception, finite-field convention, or use of the same symbol for another
object. An asymptotic claim is not automatically an effective algorithm.



Use the unit sphere $S^2\subset\R^3$. Define $d_n=\max_{x_1,\ldots,x_n\in S^2}\min_{i<j}\|x_i-x_j\|_2$. This is chordal distance; maximizing angular separation gives the same configurations because the two distances are monotonically related. Determine $d_n$ and optimal configurations in the unresolved values of $n$ specified by the source. \sref{1372}{1} \sref{1372}{2}

\paragraph{Statement recorded in the supplied source.}
For n > 14 points (except n=24), what is the maximum minimum distance between points on a unit sphere? \sref{1372}{1}

\paragraph{Residual task reported by the archive.}

The embedded review (2026-08-17) records: Solve remaining N. \sref{1372}{1}

\subsection{Short English statement}

How far apart can n points on a sphere be kept when the closest pair determines the quality of the arrangement?

\subsection{Research attempt}
\textbf{Status: UNRESOLVED.} A spherical-cap argument and a greedy construction give explicit bounds for every $n$. A Gram-matrix formulation identifies why a natural convex relaxation loses the dimension-three condition. We do not determine the optimum for a new unresolved value of $n$.

\paragraph{Context and attack.}
Musin--Tarasov [E1], checked 1 October 2026, proves the fourteen-point case by a computer-assisted contact-graph analysis. That result is outside the unresolved range requested here and is not reproduced. We first derive bounds that require no classification, then examine what a sharp certificate would have to retain. The inherited list of unresolved cases is historical; this notebook does not certify that no later isolated value has been solved.

\paragraph{Existence of an optimizer.}
The space $(S^2)^n$ is compact and the minimum of its finitely many continuous pairwise-distance functions is continuous. Hence $d_n$ is attained. A set of distinct equatorial points has positive minimum distance, so every optimizing configuration for $n\ge2$ has distinct points. Removing a point from an $(n+1)$-point configuration shows $d_{n+1}\le d_n$. Monotonicity alone does not transfer an optimal configuration between consecutive orders.

\paragraph{Disjoint-cap upper bound.}
Let $\theta$ be the minimum angular separation, so the minimum chordal separation is $\delta=2\sin(\theta/2)$. The open spherical caps of angular radius $\theta/2$ about the points are disjoint: otherwise the spherical triangle inequality would place two centres less than $\theta$ apart. A cap of angular radius $r$ has area $2\pi(1-\cos r)$, obtained by integrating the latitude bands $2\pi\sin t\,dt$. Therefore
\[
n\,2\pi(1-\cos(\theta/2))\le4\pi,
\qquad \delta\le\frac{4\sqrt{n-1}}{n}.
\]
For $n\ge2$, the lower bound $1-2/n$ on the cosine is nonnegative, so squaring it preserves the inequality used to obtain the displayed chordal bound. Boundary contacts do not affect cap area.

\paragraph{A constructive existence bound.}
Choose $\delta<2/\sqrt{n-1}$, also less than two. A spherical cap consisting of points at chordal distance less than $\delta$ from a fixed point has area exactly $\pi\delta^2$: its angular radius $\theta$ satisfies $\cos\theta=1-\delta^2/2$. Having chosen $j<n$ points, the union of their forbidden caps has area at most $j\pi\delta^2<4\pi$. There is therefore another point outside that union. Induction gives $n$ points with all distances at least $\delta$. Letting $\delta$ increase to $2/\sqrt{n-1}$ and using the compactness above yields
\[
\frac2{\sqrt{n-1}}\le d_n\le\frac{4\sqrt{n-1}}n\qquad(n\ge2).
\]
This is a proof of existence, not an efficient procedure for locating an optimal uncovered point. The constants differ by almost a factor of two, so these bounds do not solve even a single large case. A fully explicit alternative is $n$ equally spaced equatorial points, with minimum distance $2\sin(\pi/n)$; it is weaker in large $n$ but immediately supplies exact coordinates through roots of unity.

\paragraph{A second upper bound and the rank obstruction.}
For unit vectors $x_i$,
\[
0\le\left|\sum_i x_i\right|^2
=n+2\sum_{i<j}x_i\cdot x_j.
\]
Thus some pair has inner product at least $-1/(n-1)$, giving $d_n^2\le2n/(n-1)$. Equality forces all pairwise inner products to equal $-1/(n-1)$ and the vector sum to vanish. The resulting Gram matrix is $\frac n{n-1}I-\frac1{n-1}J$, with rank $n-1$. It cannot be the Gram matrix of vectors in $\mathbb R^3$ when $n>4$.

More generally, the exact optimization can be written as minimizing $\max_{i<j}G_{ij}$ over real symmetric positive semidefinite matrices with $G_{ii}=1$ and $\operatorname{rank}G\le3$. Factorization of such a matrix gives the desired unit vectors, and conversely their Gram matrix satisfies those conditions. Dropping the rank constraint admits the regular-simplex matrix above for all $n$. This proves concretely that the dimension-free semidefinite relaxation cannot by itself solve the requested spherical problem.

\paragraph{Verification and remaining gap.}
At $n=2$ both cap/greedy bounds give distance two, agreeing with antipodal points. At $n=4$ the simplex bound agrees with a regular tetrahedron; the rank obstruction starts exactly beyond that order. These analytic checks require no numerical optimization. A floating-point Gram matrix would additionally need certified positive semidefiniteness, rank control, and distance comparisons before being an exact configuration certificate.

The remaining task is a matching lower configuration and upper bound for each requested $n$. Contact graphs can encode which pairs attain the minimum, but an enumeration would have to justify its completeness and rule out all alternative realizations. The missing ingredient is that global geometric exclusion, not compactness or the ability to generate many well-spaced points.

\subsection{Sources}

{\small

\begin{itemize}
\item[E1] O. R. Musin and A. S. Tarasov, \emph{The Tammes problem for N=14}. Primary abstract checked 1 October 2026. \url{https://arxiv.org/abs/1410.2536}.


\item[\hypertarget{source:1372:1}{S1}] ULAM.AI, \emph{UnsolvedMath}, supplied \texttt{problems.json}, record 1372 (GEOM-013), statement and background. The embedded literature-review date is 2026-08-17; that is the archive’s review date, not a fresh certification. Dataset landing page: \url{https://huggingface.co/datasets/ulamai/UnsolvedMath}.

\item[\hypertarget{source:1372:2}{S2}] Tammes problem overview. \url{https://en.wikipedia.org/wiki/Tammes_problem}. Bibliographic information imported from the supplied record.

\end{itemize}

}

\label{end:1372}
\end{document}
